\begin{abstract}
Fake-news detectors score near-perfectly on their training corpus and much worse elsewhere, but cross-domain studies report this with accuracy or F1, which conflates two failures with different remedies: ranking new items badly, and ranking them well but thresholding them wrongly. We separate the two. Ten public corpora in English and Bengali (articles, headlines, claims, tweets) are audited for label direction and near-duplicates, and six detectors are evaluated train-on-one, test-on-all with five seeds, bootstrap intervals and permutation controls. The audit finds that the two most used article corpora code labels in opposite directions, which turns moderate transfer into an apparent below-chance collapse, as well as label-conflicting duplicates and tweets shared across rumour corpora. Mean ROC-AUC falls from \InAUC{} in-domain to \OffAUC{} across corpora. Decomposing the lost macro-F1 attributes \RankShareAll{} to a ranking deficit and \OpShareAll{} to the operating point. The standard label-shift correction, EM prior adaptation, enlarges the operating-point loss because calibration does not transfer (expected calibration error \ECEin{} in-domain, \ECEout{} across corpora), whereas flagging the source's base rate of target items recovers \RecRank{} of it without target labels. Twenty-two surface statistics, publisher identity or 64 training examples recover much of the in-domain signal, and frozen RoBERTa features transfer no better than TF-IDF. Pooling six training corpora raises held-out AUC from \LocoMeanSingle{} to \LocoPooled{}, while removing corpus identity from the representation lowers it, because corpus and label are entangled. Bengali clickbait detection, in contrast, survives unseen publishers and languages. We release code, per-item scores and a reporting protocol.
\end{abstract}
